problem:  
Let $\Sigma^2 \subset S^n$ be a smooth, compact, oriented immersed surface with mean curvature vector $H$ and second fundamental form $A$.  
For $\lambda>0$, define the *regularized Willmore energy*
\[
\mathcal{W}_\lambda(\Sigma) = \int_\Sigma \big(|H|^2 + \lambda |\nabla^\perp H|^2\big)\, d\mu.
\]
Assume that for each $\lambda>0$, there exists a smooth immersion $\Sigma_\lambda$ of fixed genus $g$ that is a critical point of $\mathcal{W}_\lambda$ and satisfies the scale‑invariant bound
\[
\int_{\Sigma_\lambda} |A|^2 \, d\mu \le C.
\]
Suppose further that $\Sigma_\lambda$ subconverges in the varifold sense (possibly after Möbius transformations) to a limit $\Sigma_0$ as $\lambda \to 0^+$.

**Question:**  
Is it true that any such limit $\Sigma_0$ is a smooth Willmore surface, and moreover, does there exist a *quantized curvature loss law*
\[
\lim_{\lambda\to 0^+} \mathcal{W}_\lambda(\Sigma_\lambda) = \mathcal{W}(\Sigma_0) + 4\pi k,
\]
for some integer $k\ge0$, corresponding to formation of $k$ spherical bubbles at energy scale of $4\pi$ each?  
Equivalently, do critical points of $\mathcal{W}_\lambda$ satisfy an **energy quantization** and **Γ‑limit compactness** principle analogous to that known in harmonic map and Yang–Mills theory—but now for this sixth‑order curvature functional?

Why is it a "good" problem:  
This reformulated question isolates a precise and profound phenomenon: the *energy quantization and Γ‑convergence* structure for critical points of a higher‑order Willmore‑type functional.  
It goes beyond simply asking whether the limiting surface is Willmore by proposing a conjectural **quantized energy loss mechanism**—a geometric law that potentially links the Willmore functional to the analytic framework of concentration‑compactness and gauge theory.

Several deep reasons make this a *good* problem:

1. **Conceptual simplicity** – The question is succinct: does a quantized bubble‑tree compactness hold for critical points of the “normal‑regularized” Willmore energy?  
   Its formal structure echoes classic compactness theories but transplanted into the far less explored territory of sixth‑order curvature energies.

2. **High analytical depth** – Understanding the limit involves delicate elliptic regularity and blow‑up analysis for critical points governed by a sixth‑order geometric PDE lacking conformal invariance.  
   Establishing quantization would require new monotonicity or defect measure identities.

3. **Bridging fields** – The problem links *geometric analysis* (bubbling phenomena, curvature concentration), *variational calculus* (Γ‑convergence), and *mathematical physics* (energy quantization laws known from gauge theories).  
   A resolution would reveal whether Willmore‑type energies possess an analogous quantization structure—currently unknown.

4. **Generative power** – A positive result would extend the paradigm of “quantized curvature concentration” from mean curvature or harmonic map theories to a new, higher‑order curvature setting, potentially unveiling a new geometric measure structure governing conformal surface theory in spheres.

In short, this problem is simple to state yet requires integrating geometric measure theory, analytic compactness, and high‑order PDE methods. It points directly to an unexplored frontier in the analytic geometry of surfaces, representing a genuinely *good* research-level question.